% capitolo 6: conclusioni
\chapter{Conclusioni}
\label{cap:conclusioni}

L'elaborato ha valutato se gli annunci tariffari tra Stati Uniti e Cina del biennio 2025-2026 hanno modificato la volatilità condizionata dei 
sei strumenti finanziari selezionati, con quale intensità, persistenza e con quali differenze tra asset class. L'analisi si è basata su un modello 
EGARCH(1,1) con errori t di Student, esteso con variabili \emph{dummy} definite su una finestra di due sedute per ogni evento selezionato.

Il modello minimizza il BIC per tutti e sei gli asset rispetto alle altre specificazioni considerate (sGARCH, GJR-GARCH ed EGARCH, con errori gaussiani e t di Student). 
I parametri stimati confermano i fatti stilizzati discussi nel Capitolo 2: code pesanti (da 4,2 a 6,1 gradi di libertà), persistenza elevata della volatilità ($\beta^* > 0,96$) 
e un effetto leva significativo per S\&P 500, Apple, GM e Nucor. Per l'oro l'asimmetria è invece di segno opposto, mentre per l'USD/CNH non è significativa. Poiché l'emivita degli shock va da 
circa 21 sedute per Apple a circa 86 per General Motors, l'effetto di un annuncio si prolunga per settimane oltre le due sedute della finestra.

Il test congiunto fornisce un'evidenza differenziata: per l'oro l'ipotesi nulla è rifiutata anche dopo la correzione di Holm, per Apple 
l'evidenza è intermedia (significativa al 10\%), per S\&P 500, USD/CNH e Nucor il rifiuto al 5\% vale solo prima della correzione, 
mentre per GM non emerge alcuna significatività. Il BIC preferisce sempre il modello senza \emph{dummy}, ma ciò dipende dalla penalizzazione 
dei sette parametri aggiuntivi e non è stato usato per escluderle.

Il risultato più solido riguarda il \emph{Liberation Day} (2 aprile 2025), che aumenta in modo significativo la volatilità di cinque asset 
su sei, mentre per GM l'effetto è positivo ma non significativo. Un secondo episodio rilevante è la crisi delle terre rare, con aumenti 
significativi per oro, Nucor e S\&P 500. Gli eventi di distensione si associano a riduzioni meno frequentemente significative, mentre per 
il vertice di Ginevra, il provvedimento anti-Fentanyl e la sentenza della Corte Suprema non emergono effetti stabili. I coefficienti negativi 
dell'escalation del 9 aprile non indicano un calo dell'incertezza: il modello prevede già una volatilità molto elevata dopo il \emph{Liberation Day}, e il 
coefficiente ne corregge l'effetto. Gli annunci di escalation più rilevanti (Liberation Day e terre rare) sembrano associarsi ad aumenti di volatilità più ampi e 
frequenti di quelli di distensione, e la reazione agli annunci tra gli strumenti è eterogenea: GM non mostra effetti significativi congiunti nonostante l'esposizione 
diretta agli annunci, mentre l'oro, l'S\&P 500 e Nucor reagiscono alla crisi delle terre rare in modo diverso da Apple e dal cambio. La diagnostica sui 
residui standardizzati non evidenzia criticità rilevanti.

L'analisi delle correlazioni tra i residui standardizzati mostra che i legami tra le coppie di asset analizzate sono cambiati 
in modo eterogeneo nel periodo dei dazi. La correlazione tra S\&P 500 e oro sale da 0,028 a 0,166, e la correlazione \emph{rolling} raggiunge 
il massimo dell'intero campione nel 2026: in una fase del periodo di dazi oro e azionario sembrano aver risposto allo stesso modo a shock comuni, 
con una possibile riduzione del beneficio di diversificazione dell'oro nei momenti di stress. Per Nucor e GM la correlazione scende da 0,396 a 0,220 e la 
curva \emph{rolling} si avvicina a zero, coerentemente con un impatto differenziato delle tariffe sui due settori, che secondo l'ipotesi discussa nella Sezione \ref{sec:definizione_asset} 
favorirebbero i produttori siderurgici statunitensi e penalizzerebbero GM. Per S\&P 500 e USD/CNH la correlazione media è quasi invariata, ma questo nasconde una forte 
variazione interna, in quanto la correlazione \emph{rolling} scende fino a circa $-0,7$ nel 2026, a indicare che in quella fase i ribassi dell'azionario statunitense coincidono più 
spesso con un indebolimento dello yuan.

In conclusione, lo studio mostra come gli annunci tariffari tra Stati Uniti e Cina siano associati a variazioni della volatilità 
dei mercati finanziari, in misura eterogenea tra strumenti ed eventi, e sottolinea l'importanza di modelli in grado di cogliere la 
volatilità condizionata. I risultati vanno letti come evidenza di un'associazione, ma non implicano una diretta relazione di causa-effetto: 
ogni evento è identificato da due sole osservazioni e i test sono approssimati. Come possibili sviluppi futuri, si considerano l'estensione 
dell'orizzonte temporale, l'uso di finestre di evento di ampiezza diversa e il confronto con ulteriori specificazioni, oltre all'inclusione 
di ulteriori strumenti, mercati ed eventi di natura economica e geopolitica.

%\section{Consuntivo finale}

%\section{Raggiungimento degli obiettivi}

%\section{Conoscenze acquisite}

%\section{Valutazione personale}
